\addChapter{ABSTRAK}
\vspace{0.1cm}

{
    \singlespacing
    \begin{center}
        \textbf{ABSTRAK}
    \end{center}
    \vspace{0.5cm}

    {\frenchspacing
        Proses syuting \textit{scene} film sering kali tidak dilakukan secara berurutan sesuai dengan alur cerita.
        Salah satu cara menentukan urutan syuting \textit{scene} melainkan bergantung pada efisiensi logistik, seperti biaya transportasi dan biaya \textit{talent}.
        Pemilihan urutan syuting yang kurang optimal dapat menyebabkan peningkatan biaya produksi.
        Penelitian ini membahas permasalahan pemilihan urutan syuting \textit{scene} film sebagai bentuk permasalahan \textit{Traveling Salesman Problem} (TSP) yang diselesaikan menggunakan algoritma \textit{Particle Swarm Optimization} (PSO).
        Terdapat 60 \textit{scene} yang direpresentasikan ke dalam bentuk graf lengkap sebagai \textit{vertex} dan bobot biaya syuting antar \textit{scene} sebagai \textit{edge}.
        Urutan syuting \textit{scene} dimulai dari \textit{scene} 4 dan berakhir pada \textit{scene} 4 sesuai dengan konsep TSP.
        Data diolah melalui tahap normalisasi, perhitungan fitness, dan iterasi algoritma PSO untuk memperoleh urutan \textit{scene} yang menghasilkan total biaya minimum.
        Hasil penelitian menunjukkan bahwa penerapan algoritma PSO mampu memberikan solusi urutan syuting \textit{scene} dengan total biaya minimum sebesar Rp$7.794.359,48$.
        Urutan syuting \textit{scene} yang diperoleh dari solusi yaitu 4-1C-22B-27-10-12B-32A-23B-22C-2A-34B-23A-33B-33A-19-31-1E-7A-32B-24-16B-22A-3-35B-14B-9A-28-12A-29A-21B-35A-\newline
        37-6A-25-8-7B-30A-26-16A-11-2B-19B-1D-13-5-34-6B-20-9B-15-17-16D-18-\newline
        29B-30B-16C-30C-14A-21A-36-4.
        Dengan demikian, algoritma PSO dapat menjadi alternatif solusi dalam pemilihan urutan syuting \textit{scene} pada produksi film.

        \vspace*{0.5cm}
        \noindent \textbf{Kata Kunci:} graf, \textit{particle swarm optimization} (PSO), \textit{traveling salesman problem} (TSP), urutan syuting \textit{scene}
    }
}

\pagebreak